% Job posting: https://ca.linkedin.com/jobs/view/senior-digital-fpga-engineer-at-per-vices-corporation-4465112626
\documentclass[11pt]{article}
\usepackage[left=1in,top=1in,right=1in,bottom=1in]{geometry}
\usepackage[hidelinks]{hyperref}
\usepackage{setspace}
\usepackage[T1]{fontenc}
\usepackage{newtxtext,newtxmath}
\ifdefined\pdfgentounicode
  \input{glyphtounicode}
  \pdfgentounicode=1
\fi
\setstretch{1.1}
\pagenumbering{gobble}
\begin{document}
\pagestyle{empty}

\begin{flushright}
Nishal Stanislaus Silva, PhD \\
Toronto, ON \\
nishal.silva@hotmail.com \\
+1 613 200 2030
\end{flushright}

Dear Hiring Manager,

My graduate research began with a problem that sits at the core of any digital signal chain: how to reliably extract a transient event from a noisy, continuous real-world signal. My MSc thesis at Sheffield Hallam University developed an onset-detection method using the Stockwell transform (S-transform), a time-frequency analysis technique, and was published at the IEEE Asilomar Conference on Signals, Systems, and Computers in 2018. That early grounding in time-frequency signal processing has shaped every project I have taken on since, and it is why Per Vices' work building the digital signal-processing backbone of software-defined radio hardware caught my attention.

Since then I have carried that signal-processing discipline into real-time, resource-constrained deployments. As a postdoctoral researcher at the University of Trento, I owned a full DSP-to-deployment pipeline for streaming audio and sensor data, shipping a detector that ran at 14 ms inference on a Raspberry Pi 4 (F1 0.76, 31.4\% CPU utilization, 74x faster than the 1038 ms baseline it replaced), and I designed a UDP-based edge-to-server offload architecture with characterized signal-transmission latency. I have also worked directly at the hardware boundary: at McGill University's IDMIL/CIRMMT lab I integrated an IMU over I2C with a HiFiBerry audio ADC/DAC on embedded Linux, and at MAS Holdings I interfaced computer-vision inspection systems directly with PLC-driven reject mechanisms on live manufacturing lines. All of this work has been built in C++17, with an emphasis on meeting hard latency and compute budgets rather than optimizing in the abstract.

I want to be direct about where my experience does not yet reach: I have not worked with FPGAs, HDL (VHDL/Verilog), or RF-specific hardware design, and this role's core toolchain would be new to me. What I bring instead is the signal-processing theory and the embedded-hardware integration discipline that tend to transfer quickly into that toolchain -- a track record of taking a DSP algorithm from mathematical description through to hardware-constrained, real-time deployment, and comfort working at the boundary between software and physical hardware interfaces. I would want to be evaluated on that basis, not as a substitute for FPGA/HDL experience itself.

I am a Canadian Permanent Resident based in Toronto, eligible to work without sponsorship, and available immediately. I would welcome the chance to discuss how my signal-processing and embedded-systems background could contribute to Per Vices' engineering team.

Sincerely, \\
Nishal Stanislaus Silva

\end{document}
